\begin{proofbox}

	\textbf{\textcolor{CProof}{思路提示}}

	椭圆、双曲线、抛物线的标准方程都可以从各自的定义出发，通过距离公式列出等式，再利用代数变形（平方、移项、整理）得到。变形中需注意符号处理和参数代换。

	\medskip
	\textbf{\textcolor{CProof}{正式推导}}
	\begin{description}[style=nextline, leftmargin=2.8em, labelsep=0.5em, itemsep=0.55em]
		\item[\textbf{步骤一（椭圆方程推导）：}]
			\[
				\begin{aligned}
					\sqrt{(x+c)^2+y^2}+\sqrt{(x-c)^2+y^2}&=2a \\
					\sqrt{(x+c)^2+y^2}&=2a-\sqrt{(x-c)^2+y^2} \\
					(x+c)^2+y^2&=4a^2-4a\sqrt{(x-c)^2+y^2} \\
					&\quad + (x-c)^2+y^2 \\
					4a\sqrt{(x-c)^2+y^2}&=4a^2-4cx \\
					a\sqrt{(x-c)^2+y^2}&=a^2-cx \\
					a^2[(x-c)^2+y^2]&=(a^2-cx)^2 \\
					(a^2-c^2)x^2+a^2y^2&=a^2(a^2-c^2) \\
					\frac{x^2}{a^2}+\frac{y^2}{b^2}&=1,\quad b^2=a^2-c^2.
			\end{aligned}
			\]
			\par\textbf{理由：}
			第一步用距离公式写出和等式；第二步移项使一个根号单独一边；第三步平方消去第一个根号；第四步整理出含一个根号的式子；第五步再次平方消去第二个根号；第六步展开并整理得到只含 $x^2,y^2$ 的等式；第七步用 $b^2=a^2-c^2$ 替换得到标准形式。


\item[\textbf{步骤二（双曲线方程推导）：}]
记 $r_1=\sqrt{(x+c)^2+y^2}$、$r_2=\sqrt{(x-c)^2+y^2}$，
$U=x^2+y^2+c^2$。由 $|r_1-r_2|=2a$，
\[
\begin{aligned}
r_1^2+r_2^2-2r_1r_2&=4a^2,\\
r_1r_2&=U-2a^2,\\
(r_1r_2)^2&=U^2-4c^2x^2,\\
(U-2a^2)^2&=U^2-4c^2x^2,\\
a^2(U-a^2)&=c^2x^2,\\
(c^2-a^2)x^2-a^2y^2&=a^2(c^2-a^2).
\end{aligned}
\]
令 $b^2=c^2-a^2>0$，即得 $x^2/a^2-y^2/b^2=1$。
这里引入 $r_1,r_2,U$ 只是缩短根式，不改变等式。
逆向可从标准方程得 $r_1=|cx/a+a|$、$r_2=|cx/a-a|$，
利用 $|x|\ge a$ 核验距离之差的绝对值为 $2a$。

\item[\textbf{步骤三（抛物线方程推导）：}]
			\[
				\begin{aligned}
					\sqrt{\left(x-\frac{p}{2}\right)^2+y^2}&=\left|x+\frac{p}{2}\right| \\
					\left(x-\frac{p}{2}\right)^2+y^2&=\left(x+\frac{p}{2}\right)^2 \\
					x^2-px+\frac{p^2}{4}+y^2&=x^2+px+\frac{p^2}{4} \\
					y^2&=2px.
			\end{aligned}
			\]
			\par\textbf{理由：}
			第一步根据抛物线的定义（动点到焦点的距离等于到准线的距离）列出等式；第二步两边平方消去根号和绝对值；第三步展开平方项；第四步消去 $x^2$ 和常数项，整理即得 $y^2=2px$。
	\end{description}

	\medskip
	\textbf{\textcolor{CProof}{结论回扣}}

	\textbf{椭圆、双曲线、抛物线的标准方程分别由定义通过代数变形得到。推导过程展示了“距离公式 + 平方消元”的基本方法，同时也明确了参数 $a,b,c,p$ 的几何意义。}
\end{proofbox}
